\documentclass[11pt,a4paper]{article}

% --------------------
% Packages
% --------------------
\usepackage[a4paper,margin=2.5cm]{geometry}
\usepackage{amsmath}
\usepackage{graphicx}
\usepackage[
  colorlinks=true,
  linkcolor=black,
  urlcolor=blue,
  citecolor=black
]{hyperref}
\usepackage{setspace}
\usepackage{enumitem}
\usepackage{listings}
\usepackage{xcolor}

\setstretch{1.15}

% --------------------
% Code formatting
% --------------------
\lstset{
  basicstyle=\ttfamily\small,
  breaklines=true,
  frame=single,
  numbers=left,
  numberstyle=\tiny,
  keywordstyle=\color{blue},
  commentstyle=\color{gray}
}

% --------------------
% Title
% --------------------
\title{Project Plan\\
Autonomous Pipe Inspection System}

\author{Luis Nunes and Project Team}
\date{\today}

\begin{document}

\maketitle

\tableofcontents
\newpage

% ==================================================
\section{Introduction}

This document outlines the proposed project plan for the development of an autonomous pipe inspection system. The system aims to detect, map, and classify defects within pipelines using a combination of sensing, robotics, and data processing techniques.

The project is divided into multiple subsystems, each developed by different team members. This document provides an overall structure while allowing flexibility for refinement during development.

% ==================================================
\section{Project Objectives}

The main objectives of the project are:

\begin{itemize}
  \item Design and implement a mobile pipe inspection platform
  \item Develop a reliable surface mapping system
  \item Detect structural defects such as cracks and corrosion
  \item Enable remote data transmission and control
  \item Validate system performance experimentally
\end{itemize}

% ==================================================
\section{System Overview}

The proposed system consists of the following major subsystems:

\begin{itemize}
  \item Mobile robotic platform
  \item Laser-based surface mapping system
  \item Alternative mapping sensors (stereo vision)
  \item Navigation and localisation system
  \item Communication and control interface
  \item Data processing and defect classification pipeline
\end{itemize}

These subsystems will be developed in parallel and integrated during later project stages.

% ==================================================
\section{Surface Mapping and Measurement}

\subsection{Laser Triangulation Mapping (Luis Nunes)}

\subsubsection{Overview}

This subsystem focuses on the development of a laser triangulation-based surface mapping system for internal pipe inspection.

A line laser is projected onto the internal pipe surface and observed by a calibrated camera. Using geometric triangulation, the deformation of the laser line is converted into depth information.

This method provides high-resolution local surface measurements and is suitable for detecting small-scale defects.

\subsubsection{Hardware Configuration}

The mapping system consists of:

\begin{itemize}
  \item Raspberry Pi 5
  \item Raspberry Pi Camera Module 3
  \item Red line laser module (635--660 nm)
  \item Adjustable mechanical mount
  \item Linear translation mechanism
\end{itemize}

The camera and laser are mounted at a fixed baseline and angle to enable accurate triangulation.

\subsubsection{Camera Calibration}

Camera intrinsics are obtained using a planar checkerboard calibration procedure implemented in OpenCV.

The calibration estimates:

\begin{itemize}
  \item Focal lengths
  \item Principal point
  \item Lens distortion coefficients
\end{itemize}

Calibration is performed at the final operating resolution and focus setting to ensure consistency.

\subsubsection{Laser Plane Calibration}

Laser plane calibration is used to determine the 3D equation of the projected laser sheet.

The proposed method involves:

\begin{enumerate}
  \item Capturing checkerboard images with laser enabled
  \item Estimating board pose using camera intrinsics
  \item Back-projecting laser pixels to rays
  \item Intersecting rays with the board plane
  \item Fitting a plane model
\end{enumerate}

The resulting plane parameters enable conversion from image coordinates to real-world depth.

\subsubsection{Image Processing Pipeline}

The current processing pipeline includes:

\begin{enumerate}
  \item Image acquisition
  \item Region of interest selection
  \item Gaussian smoothing
  \item HSV colour thresholding
  \item Morphological filtering
  \item Laser centreline extraction
  \item Subpixel refinement
\end{enumerate}

The pipeline outputs a set of 2D pixel coordinates representing the laser stripe.

\subsubsection{3D Reconstruction}

Each detected pixel coordinate $(u,v)$ is converted to a 3D ray:

\begin{equation}
\mathbf{r} = K^{-1}
\begin{bmatrix}
u \\ v \\ 1
\end{bmatrix}
\end{equation}

Intersection of this ray with the calibrated laser plane yields a 3D surface point.

Sequential scans are combined to form a depth profile.

\subsubsection{Alternative Mapping: Conical Mirror System (Exploratory)}

If time permits, a conical mirror combined with a point laser will be investigated.

This approach aims to project a circular laser profile onto the pipe wall, enabling full circumferential mapping in a single image.

This method will be evaluated experimentally and compared against line-based triangulation.

\subsubsection{Development Milestones}

\begin{itemize}
  \item Camera calibration and validation
  \item Stable laser extraction pipeline
  \item Laser plane calibration
  \item Depth reconstruction validation
  \item Surface accuracy assessment
  \item Integration with motion system
\end{itemize}

% ==================================================
\subsection{Stereo Vision Mapping (To Be Completed)}

This section will describe the development of a stereo vision-based surface reconstruction system.

Details to be completed by responsible team members.

% ==================================================
\section{Robotic Platform Design (To Be Completed)}

This section will outline the mechanical and actuation design of the mobile inspection robot.

Topics include:

\begin{itemize}
  \item Locomotion mechanism
  \item Structural design
  \item Power distribution
  \item Actuation systems
\end{itemize}

Details to be added by the robotics team.

% ==================================================
\section{Navigation and Localisation (To Be Completed)}

This section will describe methods for estimating robot position and orientation within pipes.

Potential sensors include:

\begin{itemize}
  \item Ultrasonic sensors
  \item LiDAR
  \item Inertial sensors
\end{itemize}

Details to be added by responsible team members.

% ==================================================
\section{Communication and Control (To Be Completed)}

This section will outline remote communication and system control methods.

Topics include:

\begin{itemize}
  \item Wireless communication
  \item Data transmission
  \item Remote monitoring
  \item Control interfaces
\end{itemize}

% ==================================================
\section{Data Processing and Defect Classification (To Be Completed)}

This section will describe post-processing and machine learning methods for defect identification.

Potential techniques include:

\begin{itemize}
  \item Feature extraction
  \item Statistical analysis
  \item Machine learning classifiers
  \item Deep learning models
\end{itemize}

% ==================================================
\section{System Integration and Testing}

System integration will be performed in stages:

\begin{enumerate}
  \item Subsystem validation
  \item Mechanical-electrical integration
  \item Sensor synchronisation
  \item Software integration
  \item End-to-end testing
\end{enumerate}

Testing will include controlled laboratory experiments and representative pipe environments.

% ==================================================
\section{Project Management and Timeline}

Each subsystem will follow an internal development schedule.

Key milestones will be mapped to a consolidated project Gantt chart at a later stage once subsystem development stabilises.

This approach allows adaptive planning while maintaining overall project coherence.

% ==================================================
\section{Risk Assessment}

Key project risks include:

\begin{itemize}
  \item Insufficient mapping accuracy
  \item Mechanical instability
  \item Communication failures
  \item Integration delays
  \item Limited testing time
\end{itemize}

Mitigation strategies will be developed throughout the project.

% ==================================================
\section{Conclusion}

This project plan provides a structured framework for the development of an autonomous pipe inspection system.

The laser triangulation subsystem forms a core component of the surface mapping capability and will be developed and validated early in the project.

The document will be updated iteratively as development progresses.

\end{document}
